\section{Appendix A: Formal Notation, Assumptions, Theorem Register, and Proof Map}
\label{sec:prof-app-a}

This appendix is the compact formal reference for the professor-facing IEEE draft.  It
collects the notation, the standing assumptions, the theorem register, and the proof-map
framing used in the main theory bridge.  The goal is to keep the body readable while
making every strong claim traceable to a formal object or a local source file.

\subsection{Notation}
\begin{table}[t]
\centering
\caption{Canonical notation for the ORIUS professor draft.}
\label{tab:prof-notation}
\begin{tabular}{p{0.26\textwidth}p{0.66\textwidth}}
\toprule
\textbf{Symbol / Term} & \textbf{Meaning} \\
\midrule
\(x_t\) & True plant state. \\
\(y_t\) & Telemetry or observation at time \(t\). \\
\(\mathcal{H}_t\) & Observation history generated by \(y_{0:t}\). \\
\(w_t\in[0,1]\) & Reliability score extracted from telemetry history. \\
\(U_t\) & Observation-consistent state set. \\
\(\mathcal{S}\) & Safe state set defined by the domain safety predicate. \\
\(\mathcal{A}_t^{\mathrm{safe}}(U_t)\) & Tightened safe-action set under uncertainty set \(U_t\). \\
\(a_t^{\mathrm{cand}}\) & Nominal candidate action produced by the upstream controller. \\
\(a_t^{\mathrm{safe}}\) & Repaired or fallback action emitted by ORIUS. \\
\(V_T\) & Episode-level true-state violation count. \\
\(\alpha\) & Miscoverage or per-step risk budget. \\
\(\bar{w}\) & Mean reliability across the horizon. \\
\(\phi(x)\) & Safety predicate defining the domain safe set. \\
\bottomrule
\end{tabular}
\end{table}

\subsection{Standing assumptions}
The formal objects used in the body are grounded in the assumption register from the
local thesis sources.  The most important assumptions are:
\begin{itemize}
  \item \textbf{A1: bounded model error.} The nominal dynamics error is bounded.
  \item \textbf{A2: bounded telemetry error.} Observation error is controlled by a
  monotone function of reliability.
  \item \textbf{A3: feasible safe repair.} If a safe state exists, the tightened action
  set is non-empty.
  \item \textbf{A4: known dynamics.} The local transition law is available to the
  controller or safety layer.
  \item \textbf{A5: monotone bounded inflation.} Lower reliability never shrinks the
  uncertainty surface.
  \item \textbf{A6: bounded detector lag.} Degradation is detected within a finite lag.
  \item \textbf{A7: causal certificate rule.} Certificates depend only on current and
  past information.
  \item \textbf{A8: admissible fallback.} A safe fallback policy exists under the
  admissible fault model.
\end{itemize}
These assumptions are not decorative.  They are the minimum conditions under which the
theory bridge can claim that the universal runtime contract is not just readable, but
sound.

\subsection{Theorem register}
\begin{table}[t]
\centering
\caption{Professor-facing theorem register for the ORIUS draft.}
\label{tab:prof-theorem-register}
\begin{tabular}{p{0.12\textwidth}p{0.34\textwidth}p{0.42\textwidth}}
\toprule
\textbf{Tag} & \textbf{Claim} & \textbf{Role in the paper} \\
\midrule
T9 & Reliability-aware repair is structurally necessary. & Establishes that ignoring \(w_t\) yields linear violation scaling under persistent degradation. \\
T10 & Degraded observation induces a reliability-risk frontier. & Lower bound showing that the risk envelope cannot be removed by controller choice alone. \\
T11 & Structural transfer through the universal adapter contract. & Shows that the battery proof pattern transfers when a domain supplies the required runtime objects. \\
\bottomrule
\end{tabular}
\end{table}

% tbl_theorem_dependency.tex
% Theorem dependency and proof completeness (Applied Math Prof critique).
\begin{table}[htbp]
\centering
\caption{DC3S Theorem Dependency and Proof Completeness. Each formal result is implemented in code and tested by at least one automated test.}
\label{tbl:theorem_dependency}
\resizebox{\textwidth}{!}{%
\begin{tabular}{lllll}
\toprule
\textbf{Result} & \textbf{Name} & \textbf{Depends On} & \textbf{Proof Method} & \textbf{Code Reference} \\
\midrule
Defn.~1 & OQE reliability weight $w_t$ & — & Definition & \texttt{quality.py:compute\_reliability()} \\
Prop.~1 & Inflation monotonicity & Defn.~1 & Analytic & \texttt{coverage\_theorem.py:verify\_inflation\_geq\_one()} \\
Thm.~1  & Marginal coverage guarantee & Prop.~1, Defn.~1 & Conformal CP theory & \texttt{coverage\_theorem.py:assert\_coverage\_guarantee()} \\
Thm.~3  & Expected violation bound & Thm.~1 & Expectation calculus & \texttt{coverage\_theorem.py:compute\_expected\_violation\_bound()} \\
Thm.~9  & Group-conditional coverage & Thm.~1 & Mondrian CP & \texttt{coverage\_theorem.py:mondrian\_group\_coverage()} \\
Thm.~10 & Hoeffding high-prob.\ bound & Thm.~3 & Hoeffding inequality & \texttt{coverage\_theorem.py:hoeffding\_violation\_bound()} \\
Thm.~11 & Byzantine OQE reliability bound & Defn.~1 & Trimmed-mean theory & \texttt{quality.py:compute\_reliability\_robust()} \\
\bottomrule
\multicolumn{5}{l}{\small All results are unit-tested in \texttt{tests/test\_conditional\_coverage.py} and \texttt{tests/test\_adversarial\_faults.py}.}\\
\end{tabular}%
}
\end{table}


\subsection{Proof-map framing}
The proof map used in the main text has three levels.
\begin{enumerate}
  \item \textbf{Reference-domain proofs.}  The battery ladder and the supporting
  theorem surface establish the reference proof objects.  These are the source of the
  concrete runtime semantics and the strongest empirical closure.
  \item \textbf{Transfer proofs.}  The universality chapter isolates the objects that a
  new domain must supply to inherit the same one-step safety statement.
  \item \textbf{Promotion proofs.}  The cross-domain synthesis and deployment sections
  use the universality gate to decide whether a domain may be treated as defended,
  bounded, gated, or experimental in the article body.
\end{enumerate}

In practice, the proof map is read as a dependency chain:
\[
  \text{A1--A13} \Longrightarrow
  \text{battery reference ladder} \Longrightarrow
  \text{T9/T10/T11} \Longrightarrow
  \text{universality gate} \Longrightarrow
  \text{domain promotion status}.
\]
This is the simplest way to keep the professor-facing draft universal-first without
collapsing evidence tiers.  The theorem register tells the reader what is proved, and the
proof map tells the reader what has to be true before those proofs can be reused.

\subsection{How to use this appendix}
When a claim in the main IEEE text depends on a formal step, cite the theorem tag from
Table~\ref{tab:prof-theorem-register} and the assumption ID from the register above.
When a claim depends on a battery-specific derivation, use the proof-map framing to make
clear whether the claim is a reference-domain result, a transfer theorem, or a
promotion-level statement.  This avoids the common failure mode in stitched monographs:
letting broad architectural language outrun the actual proof obligations.
